% =============================================================
%  00_cheatsheet.tex — Common Lisp Cheat Sheet
%  1-2 pagine dense, due colonne, link interni alla reference
% =============================================================
\chapter*{Cheat Sheet}
\addcontentsline{toc}{chapter}{Cheat Sheet}
\markboth{Cheat Sheet}{Cheat Sheet}

{\small
\begin{multicols}{2}

% ---- Aritmetica --------------------------------------------
\noindent\textbf{\hyperref[sec:numeri]{Aritmetica}}\\
\begin{tabular}{@{}ll@{}}
\fn{(+ a b \ldots)}    & somma \\
\fn{(- a b)}           & differenza \\
\fn{(* a b \ldots)}    & prodotto \\
\fn{(/ a b)}           & quoziente (esatto) \\
\fn{(mod a b)}         & resto (segno di \texttt{b}) \\
\fn{(rem a b)}         & resto (segno di \texttt{a}) \\
\fn{(expt b e)}        & $b^e$ \\
\fn{(sqrt x)}          & radice quadrata \\
\fn{(abs x)}           & valore assoluto \\
\fn{(1+ x)} / \fn{(1- x)} & inc/dec \\
\fn{(floor x)}         & parte intera $\lfloor x \rfloor$ \\
\fn{(ceiling x)}       & $\lceil x \rceil$ \\
\fn{(round x)}         & arrotonda \\
\fn{(max a b)} / \fn{(min a b)} & massimo/minimo \\
\end{tabular}

\medskip
% ---- Confronti --------------------------------------------
\noindent\textbf{\hyperref[sec:numeri]{Confronti numerici}}\\
\begin{tabular}{@{}ll@{}}
\fn{(= a b)}  & uguale \\
\fn{(/= a b)} & diverso \\
\fn{(< a b)}  \fn{(> a b)} \fn{(<= a b)} \fn{(>= a b)} & ordine \\
\end{tabular}

\medskip
% ---- Booleani ---------------------------------------------
\noindent\textbf{\hyperref[sec:booleani]{Booleani}}\\
\begin{tabular}{@{}ll@{}}
\fn{T} / \fn{NIL}    & vero / falso (e lista vuota) \\
\fn{(and a b \ldots)} & e logico (short-circuit) \\
\fn{(or a b \ldots)}  & o logico (short-circuit) \\
\fn{(not x)}          & negazione \\
\end{tabular}

\medskip
% ---- Variabili e binding ----------------------------------
\noindent\textbf{\hyperref[sec:variabili]{Definizioni e binding}}\\
\begin{tabular}{@{}ll@{}}
\fn{(defun nome (p\ldots) corpo)} & funzione \\
\fn{(defparameter *x* val)}       & variabile globale \\
\fn{(defvar *x* val)}             & globale (no-overwrite) \\
\fn{(defconstant +k+ val)}        & costante \\
\fn{(let ((x v)\ldots) corpo)}    & binding locale \\
\fn{(let* ((x v)\ldots) corpo)}   & binding sequenziale \\
\fn{(setf x val)}                 & assegnamento \\
\end{tabular}

\medskip
% ---- Condizioni -------------------------------------------
\noindent\textbf{\hyperref[sec:condizioni]{Condizioni}}\\
\begin{tabular}{@{}ll@{}}
\fn{(if test then else)}   & ramo singolo \\
\fn{(cond (t1 e1)\ldots)}  & rami multipli \\
\fn{(when test corpo\ldots)}   & solo se vero \\
\fn{(unless test corpo\ldots)} & solo se falso \\
\fn{(case val (k e)\ldots)}    & switch su valore \\
\end{tabular}

\columnbreak

% ---- Liste ------------------------------------------------
\noindent\textbf{\hyperref[sec:liste]{Liste — costruzione e accesso}}\\
\begin{tabular}{@{}ll@{}}
\fn{(cons x lst)}        & prependi \\
\fn{(list a b \ldots)}   & crea lista \\
\fn{(append l1 l2)}      & concatena \\
\fn{(reverse lst)}       & inverti \\
\fn{(car lst)} / \fn{(first lst)}  & primo \\
\fn{(cdr lst)} / \fn{(rest lst)}   & resto \\
\fn{(cadr lst)} / \fn{(second lst)}& secondo \\
\fn{(nth n lst)}         & n-esimo (0-based) \\
\fn{(last lst)}          & ultima cons-cell \\
\fn{(length lst)}        & lunghezza \\
\fn{(null lst)}          & lista vuota? \\
\fn{(member x lst)}      & appartiene? \\
\end{tabular}

\medskip
% ---- Higher-order -----------------------------------------
\noindent\textbf{\hyperref[sec:higherorder]{Higher-order}}\\
\begin{tabular}{@{}ll@{}}
\fn{(mapcar \#'f lst)}       & trasforma \\
\fn{(mapcan \#'f lst)}       & trasforma+appiattisce \\
\fn{(remove-if \#'p lst)}    & rimuovi se vero \\
\fn{(remove-if-not \#'p lst)}& tieni se vero \\
\fn{(find-if \#'p lst)}      & primo match \\
\fn{(every \#'p lst)}        & tutti veri? \\
\fn{(some \#'p lst)}         & almeno uno vero? \\
\fn{(reduce \#'f lst)}       & fold sinistro \\
\fn{(sort lst \#'<)}         & ordina (distruttivo) \\
\fn{(count-if \#'p lst)}     & conta match \\
\fn{(funcall f a b)}         & chiama f \\
\fn{(apply f args)}          & chiama f con lista args \\
\fn{\#'nome}                 & oggetto funzione \\
\fn{(lambda (x) corpo)}      & funzione anonima \\
\end{tabular}

\medskip
% ---- Iterazione -------------------------------------------
\noindent\textbf{\hyperref[sec:iterazione]{Iterazione}}\\
\begin{tabular}{@{}ll@{}}
\fn{(dotimes (i n) corpo)}         & 0..n-1 \\
\fn{(dolist (x lst) corpo)}        & su lista \\
\fn{(loop for x in lst do \ldots)} & loop su lista \\
\fn{(loop for i from a to b do \ldots)} & loop numerico \\
\fn{(loop \ldots collect e)}       & raccogli risultati \\
\fn{(loop \ldots sum e)}           & somma \\
\end{tabular}

\medskip
% ---- Strutture dati specializzate -------------------------
\noindent\textbf{\hyperref[sec:hashtable]{Hash table}}\\
\begin{tabular}{@{}ll@{}}
\fn{(make-hash-table)}              & crea \\
\fn{(gethash k ht)}                 & leggi \\
\fn{(setf (gethash k ht) v)}        & scrivi \\
\fn{(remhash k ht)}                 & elimina \\
\fn{(maphash \#'f ht)}              & itera \\
\end{tabular}

\medskip
% ---- Array ------------------------------------------------
\noindent\textbf{\hyperref[sec:array]{Array e vettori}}\\
\begin{tabular}{@{}ll@{}}
\fn{\#(1 2 3)}                    & vettore letterale \\
\fn{(make-array n)}               & crea vettore \\
\fn{(make-array '(r c))}          & crea matrice \\
\fn{(aref a i)} / \fn{(aref m r c)} & accesso \\
\fn{(setf (aref a i) v)}          & assegna \\
\fn{(length a)}                   & dimensione \\
\end{tabular}

\medskip
% ---- Defstruct --------------------------------------------
\noindent\textbf{\hyperref[sec:strutture]{Strutture (defstruct)}}\\
\begin{tabular}{@{}ll@{}}
\fn{(defstruct nome slot\ldots)}       & definisci \\
\fn{(make-nome :slot val)}             & crea istanza \\
\fn{(nome-slot ist)}                   & accedi \\
\fn{(setf (nome-slot ist) v)}          & modifica \\
\end{tabular}

\medskip
% ---- I/O --------------------------------------------------
\noindent\textbf{\hyperref[sec:io]{I/O}}\\
\begin{tabular}{@{}ll@{}}
\fn{(read)}                    & leggi da stdin \\
\fn{(read-line)}               & leggi riga \\
\fn{(print x)}                 & stampa + newline \\
\fn{(princ x)}                 & stampa senza quote \\
\fn{(format t "\textasciitilde A\textasciitilde \%"\ x)} & formattato \\
\fn{\textasciitilde A} \fn{\textasciitilde D} \fn{\textasciitilde F} \fn{\textasciitilde \%} & gen/int/float/newline \\
\end{tabular}

\medskip
% ---- Macro ------------------------------------------------
\noindent\textbf{\hyperref[sec:macro]{Macro}}\\
\begin{tabular}{@{}ll@{}}
\fn{(defmacro nome (p\ldots) corpo)} & definisci \\
\fn{`(lista ,expr ,@lista)}          & backquote/unquote \\
\fn{(gensym)}                        & simbolo unico \\
\end{tabular}

\medskip
% ---- Predicati --------------------------------------------
\noindent\textbf{\hyperref[sec:predicati]{Predicati comuni}}\\
\begin{tabular}{@{}ll@{}}
\fn{(null x)}      & NIL? \\
\fn{(atom x)}      & atomo (non cons)? \\
\fn{(listp x)}     & lista? \\
\fn{(numberp x)}   & numero? \\
\fn{(symbolp x)}   & simbolo? \\
\fn{(stringp x)}   & stringa? \\
\fn{(eq a b)}      & identità oggetto \\
\fn{(eql a b)}     & identità + numeri/char \\
\fn{(equal a b)}   & uguaglianza strutturale \\
\end{tabular}

\end{multicols}

\vfill
\noindent\rule{\textwidth}{0.4pt}\\[2pt]
{\footnotesize
\textbf{Legenda:}
\essential\ = esercizio essenziale\quad
\optional\ = esercizio collaterale\quad
\unaStella\ facile\quad
\dueStelle\ medio\quad
\treStelle\ difficile\quad
\quattroStelle\ avanzato\quad
\cinqueStelle\ progetto
}
} % fine \small
